\documentclass{homework}

% The following commands set up the material that appears
% in the header.
\doclabel{Math F404: Homework 1}
\docauthor{Christopher Munoz}
\docdate{September 1, 2026}

% Some common macros
\newcommand{\Reals}{\ensuremath{\mathbb R}}% Gives you a shortcut for writing the blackboard R for the real numbers - \RR
\newcommand{\Nats}{\ensuremath{\mathbb N}} % Gives you a shortcut for writing the blackboard N for the natural numbers - \NN
\newcommand{\Ints}{\ensuremath{\mathbb Z}} % Gives you a shortcut for writing the blackboard Z for the integer numbers - \ZZ

\begin{document}
\begin{problems}

\problem Sutherland 2.1 :
Suppose that $C, D$ are subsets of a set $X$. Prove that
\[
(X \setminus C) \cap D = D \setminus C.
\]
\solution
  \begin{proof}
$ $\newline In order to prove equality among sets, we must show that they are subsets of each other. 
    \textbf{Forward Case:} Let $x \in [(X \setminus C) \cap D]$. Unfolding the definitions:
    \begin{align*}
      x \in [( X \setminus C) \cap D] & \implies x \in (X \setminus C) \land x \in D \\
      & \implies (x \in X \land x \notin C) \land x \in D \\
      & \implies  x \in X \land x \notin C \land x \in D \\
      & \implies x \in D \land x \notin C && \text{$C \subseteq X$ and $D \subseteq X$}\\
      & \implies x \in [D \setminus C]
    \end{align*}
    Thus $(X \setminus C) \cap D \subseteq D \setminus C$. \newline
    \textbf{Backwards case}: Let $x \in [D \setminus C]$. Unfolding the definitions:
      \begin{align*}
        x \in [D \setminus C] &\implies x \in D \land x \notin C \\
        &\implies x \in D \land x \in X \land x \notin C && \text{$D \subseteq X \implies x \in D \land x \in X$}\\
        &\implies (x \in X \land x \notin C ) \land x \in D \\
        &\implies (X \setminus C) \cap D
      \end{align*}
      Thus $D \setminus X \subseteq (x \setminus C) \cap D$.
      From these two cases it follows that $(X \setminus C) \cap D = D \setminus C$.
\end{proof}

\problem Sutherland 2.2 :
Suppose that $A, V$ are subsets of a set $X$. Prove that
\[
A \setminus (V \cap A) = A \cap (X \setminus V).
\]
\solution
  \begin{proof}
$ $\newline In order to prove equality among sets, we must show that they are subsets of each other.
    \textbf{Forward Case:} Let $x \in [A \setminus (V \cap A)]$. Unfolding the definitions:
    \begin{align*}
      x \in [A \setminus (V \cap A)] &\implies x \in A \land x \notin (V \cap A) \\
      &\implies x \in A \land \neg(x \in V \land x \in A) \\
      &\implies x \in A \land (x \notin V \lor x \notin A) \\
      &\implies x \in A \land x \notin V && \text{since $x \in A$ rules out $x \notin A$}\\
      &\implies x \in A \land x \in X \land x \notin V && \text{$A \subseteq X$}\\
      &\implies x \in A \land (x \in X \land x \notin V) \\
      &\implies x \in A \land x \in [X \setminus V] \\
      &\implies x \in [A \cap (X \setminus V)]
    \end{align*}
    Thus $A \setminus (V \cap A) \subseteq A \cap (X \setminus V)$. \newline
    \textbf{Backwards Case:} Let $x \in [A \cap (X \setminus V)]$. Unfolding the definitions:
      \begin{align*}
        x \in [A \cap (X \setminus V)] &\implies x \in A \land x \in [X \setminus V] \\
        &\implies x \in A \land x \in X \land x \notin V \\
        &\implies x \in A \land x \notin V \\
        &\implies x \in A \land \neg(x \in V \land x \in A) && \text{since $x\notin V$ alone kills the conjunction}\\
        &\implies x \in A \land x \notin (V \cap A) \\
        &\implies x \in [A \setminus (V \cap A)]
      \end{align*}
      Thus $A \cap (X \setminus V) \subseteq A \setminus (V \cap A)$.
      From these two cases it follows that $A \setminus (V \cap A) = A \cap (X \setminus V)$.
\end{proof}

\problem Sutherland 2.5 \textbf{AI-free}:
Suppose that $U_1, U_2$ are subsets of a set $X$ and that $V_1, V_2$ are
subsets of a set $Y$. Prove that
\[
(U_1 \times V_1) \cap (U_2 \times V_2) = (U_1 \cap U_2) \times (V_1 \cap V_2).
\]
\solution
  \begin{proof}
$ $\newline Since both sides are subsets of $X \times Y$, we argue by unfolding definitions for an arbitrary pair $(x,y)$, note that the hypothesis is flavor text.
    \begin{align*}
      (x,y) \in [(U_1 \times V_1) \cap (U_2 \times V_2)]
      &\iff (x,y) \in (U_1 \times V_1) \land (x,y) \in (U_2 \times V_2) \\
      &\iff (x \in U_1 \land y \in V_1) \land (x \in U_2 \land y \in V_2) \\
      &\iff (x \in U_1 \land x \in U_2) \land (y \in V_1 \land y \in V_2) && \text{rearranging conjuncts}\\
      &\iff x \in (U_1 \cap U_2) \land y \in (V_1 \cap V_2) \\
      &\iff (x,y) \in [(U_1 \cap U_2) \times (V_1 \cap V_2)]
    \end{align*}
    Since $(x,y)$ was arbitrary, $(U_1 \times V_1) \cap (U_2 \times V_2) = (U_1 \cap U_2) \times (V_1 \cap V_2)$.
\end{proof}

\problem Sutherland 2.7 as follows:
Suppose $\sim$ is an equivalence relation on a set $X$. Show that the equivalence classes partition
$X$ in the following sense. Let $\mathcal{C}$ denote the collection of equivalence classes. Then
$\bigcup_{C \in \mathcal{C}} C = X$, and if $C, C' \in \mathcal{C}$ then either $C = C'$ or
$C \cap C' = \emptyset$.
\solution
  \begin{proof}
$ $\newline For $x \in X$, write $[x] = \{y \in X : y \sim x\}$ for the equivalence class of $x$; note $\mathcal{C} = \{[x] : x \in X\}$. \newline
    \textbf{Part 1 (the union is $X$):}
    Each $[x] \subseteq X$ by definition, so $\bigcup_{C \in \mathcal{C}} C \subseteq X$. Conversely, for any $x \in X$, reflexivity gives $x \sim x$, so $x \in [x]$ and $[x] \in \mathcal{C}$; hence $x \in \bigcup_{C \in \mathcal{C}} C$. Thus $X \subseteq \bigcup_{C \in \mathcal{C}} C$, and so $\bigcup_{C \in \mathcal{C}} C = X$. \newline
    \textbf{Part 2 (equal or disjoint):} Suppose $C, C' \in \mathcal{C}$ with $C \cap C' \neq \emptyset$; say $C = [a]$, $C' = [b]$. We show $C = C'$.
    \begin{align*}
      z \in C \cap C' &\implies z \sim a \land z \sim b \\
      &\implies a \sim z \land z \sim b && \text{symmetry}\\
      &\implies a \sim b && \text{transitivity}
    \end{align*}
    Now let $w \in C = [a]$, so $w \sim a$. Since $a \sim b$, transitivity gives $w \sim b$, so $w \in [b] = C'$. Thus $C \subseteq C'$. By symmetry $b \sim a$, and the identical argument with the roles of $a,b$ reversed gives $C' \subseteq C$. Hence $C = C'$.

    So for any $C, C' \in \mathcal{C}$, either $C \cap C' = \emptyset$ or $C = C'$.
\end{proof}

\problem Give examples of collections of open intervals $I_i = (a_i, b_i)$ with $i \in \Nats$ such that
\begin{subproblems}
\item $\bigcap_{i \in \Nats} I_i = [0,1]$
\item $\bigcap_{i \in \Nats} I_i = (0,1)$
\item $\bigcap_{i \in \Nats} I_i = (0,1]$
\end{subproblems}
For part (a) give a formal proof. For the remainder, just give an example and a brief
justification.
\subsol
  \textbf{Example:} $I_i = \left(-\dfrac{1}{i}, 1 + \dfrac{1}{i}\right)$ for $i \in \Nats$.
  \begin{proof}
$ $\newline \textbf{Forward Case ($\subseteq$):} Let $x \in \bigcap_{i\in\Nats} I_i$. Then for every $i$,
    \[
    -\frac1i < x < 1+\frac1i.
    \]
    Since $x > -\frac1i$ for every $i$, $x$ is an upper bound for $\left\{-\frac1i : i \in \Nats\right\}$, whose supremum is $0$. Hence
    \[
    x \ge \sup\left\{-\tfrac1i\right\} = 0.
    \]
    Since $x < 1+\frac1i$ for every $i$, $x$ is a lower bound for $\left\{1+\frac1i : i \in \Nats\right\}$, whose infimum is $1$. Hence
    \[
    x \le \inf\left\{1+\tfrac1i\right\} = 1.
    \]
    So $x \in [0,1]$, giving $\bigcap_{i\in\Nats} I_i \subseteq [0,1]$. \newline
    \textbf{Backwards Case ($\supseteq$):} Let $x \in [0,1]$. For every $i \in \Nats$,
    \[
    -\frac1i < 0 \le x \le 1 < 1+\frac1i,
    \]
    so $x \in I_i$. Since $i$ was arbitrary, $x \in \bigcap_{i\in\Nats} I_i$, giving $[0,1] \subseteq \bigcap_{i\in\Nats} I_i$.
 
    From these two cases, $\bigcap_{i \in \Nats} I_i = [0,1]$.
  \end{proof}


\subsol
  \textbf{Example:} $I_i = (0,1)$ for every $i \in \Nats$. \newline
  \textbf{Justification:} The intersection of countably infinite copies of the same set is just that set.

\subsol
  \textbf{Example:} $I_i = \left(0,\, 1+\dfrac1i\right)$ for $i \in \Nats$. \newline
  \textbf{Justification:} The lower endpoint is fixed at $0$, so $x \in I_i$ for all $i$ forces $x > 0$ regardless of $i$. The upper endpoints $1+\frac1i$ will converge to 1.
\problem Sutherland 3.2 (No formal proof needed):
Let $f : \Reals \to \Reals$ be defined by $f(x) = \sin x$. Describe the sets:
\[
f([0, \pi/2]), \quad f([0,\infty)), \quad f^{-1}([0,1]), \quad f^{-1}([0,1/2]), \quad f^{-1}([-1,1]).
\]
\solution
$ $\newline Since $\sin$ increases from $0$ to $1$ on $[0,\pi/2]$:
\[
f([0,\pi/2]) = [0,1].
\]
Since $\sin$ attains every value in $[-1,1]$already on, e.g., $[0,2\pi] \subseteq [0,\infty)$:
\[
f([0,\infty)) = [-1,1].
\]
$\sin x \in [0,1]$ exactly when $\sin x \ge 0$, i.e.\ on each interval $[2k\pi, (2k+1)\pi]$, $k \in \Ints$:
\[
f^{-1}([0,1]) = \bigcup_{k \in \Ints} [2k\pi,\ (2k+1)\pi].
\]
$\sin x \in [0,1/2]$ on each ``hump'' of $\sin$, restricted near its base (below the points where $\sin x = 1/2$, namely $x = \pi/6, 5\pi/6$ up to shifts of $2\pi$):
\[
f^{-1}([0,1/2]) = \bigcup_{k \in \Ints} \left( \left[2k\pi,\ 2k\pi+\tfrac{\pi}{6}\right] \cup \left[2k\pi+\tfrac{5\pi}{6},\ 2k\pi+\pi\right] \right).
\]
Since $\sin x \in [-1,1]$ for every real $x$:
\[
f^{-1}([-1,1]) = \Reals.
\]

\problem Sutherland 3.3 :
Suppose that $f : X \to Y$ and $g : Y \to Z$ are maps and $U \subseteq Z$.
Prove that
\[
(g \circ f)^{-1}(U) = f^{-1}(g^{-1}(U)).
\]
\solution
  \begin{proof}
$ $\newline For arbitrary $x \in X$, every step below is an $\iff$, so no case split is needed:
    \begin{align*}
      x \in [(g\circ f)^{-1}(U)] &\iff (g \circ f)(x) \in U \\
      &\iff g(f(x)) \in U \\
      &\iff f(x) \in [g^{-1}(U)] \\
      &\iff x \in [f^{-1}(g^{-1}(U))]
    \end{align*}
    Since $x$ was arbitrary, $(g\circ f)^{-1}(U) = f^{-1}(g^{-1}(U))$.
\end{proof}

\problem Sutherland 3.4
Let $f : \Reals \to \Reals^2$ be defined by $f(x) = (x, 2x)$. Describe the sets:
\[
f([0,1]), \quad f^{-1}([0,1] \times [0,1]), \quad f^{-1}(D)
\]
where $D = \{(x,y) \in \Reals^2 : x^2 + y^2 \le 1\}$.
\solution
$ $\newline \textbf{For} $\mathbf{f([0,1])} = \{(x,2x) : x \in [0,1]\}$, the line segment from $(0,0)$ to $(1,2)$.

\textbf{For} $\mathbf{f^{-1}([0,1]\times[0,1])}$: $f(x) = (x,2x) \in [0,1]\times[0,1]$ iff $0 \le x \le 1$ and $0 \le 2x \le 1$, i.e.\ $0 \le x \le \tfrac12$. So
\[
f^{-1}([0,1]\times[0,1]) = \left[0, \tfrac12\right].
\]

\textbf{For} $\mathbf{f^{-1}(D)}$: $f(x) = (x,2x) \in D$ iff $x^2 + (2x)^2 \le 1$, i.e.\ $5x^2 \le 1$, i.e.\ $|x| \le \tfrac{1}{\sqrt5}$. So
\[
f^{-1}(D) = \left[-\tfrac{1}{\sqrt5},\ \tfrac{1}{\sqrt5}\right].
\]

\problem Sutherland 3.8 (Make your proof as short as possible): 
Let $f : X \to Y$ be a map and let $A, B$ be subsets of $X$. Prove
that $f(A \setminus B) = f(A) \setminus f(B)$ if and only if $f(A \setminus B) \cap f(B) = \emptyset$.
Deduce that if $f$ is injective then $f(A \setminus B) = f(A) \setminus f(B)$.
\solution
  \begin{proof}
$ $\newline \textbf{Forward Case ($\implies$):} Suppose $f(A\setminus B) = f(A)\setminus f(B)$. By definition, $f(A)\setminus f(B)$ is disjoint from $f(B)$, so
    \[
    f(A \setminus B) \cap f(B) = \left[f(A)\setminus f(B)\right] \cap f(B) = \emptyset.
    \]
    \textbf{Backwards Case ($\impliedby$):} Suppose $f(A\setminus B) \cap f(B) = \emptyset$.
    \begin{itemize}
      \item[$\subseteq$:] Let $y \in f(A\setminus B)$. Since $A\setminus B \subseteq A$, $y \in f(A)$. Since $f(A\setminus B) \cap f(B) = \emptyset$, $y \notin f(B)$. So $y \in f(A)\setminus f(B)$.
      \item[$\supseteq$:] Let $y \in f(A)\setminus f(B)$, so $y = f(x)$ for some $x \in A$, and $y \notin f(B)$. If $x \in B$ then $y = f(x) \in f(B)$, contradicting $y \notin f(B)$; so $x \notin B$. Thus $x \in A\setminus B$, so $y \in f(A\setminus B)$. (This inclusion needs no hypothesis.)
    \end{itemize}
    Hence $f(A\setminus B) = f(A)\setminus f(B)$.

    \textbf{Deduction:} Suppose $f$ is injective. If $y \in f(A\setminus B) \cap f(B)$, then $y = f(x_1)$ for some $x_1 \in A\setminus B$ and $y = f(x_2)$ for some $x_2 \in B$. Injectivity forces $x_1 = x_2$, so $x_1 \in B$ — contradicting $x_1 \in A \setminus B$. Hence $f(A\setminus B)\cap f(B) = \emptyset$, and by the equivalence above, $f(A\setminus B) = f(A)\setminus f(B)$.
\end{proof}

\end{problems}
\end{document}
